\documentclass{article}

\usepackage[utf8]{inputenc}
\usepackage[T2A]{fontenc}
\usepackage[russian]{babel}
\usepackage{amsmath,amssymb}
\usepackage{amsthm}
\usepackage{geometry}
\geometry{left=2.5cm,right=2.5cm,top=2.5cm,bottom=2.5cm}
\usepackage[unicode]{hyperref}

\title{Упрощение перебора минимальных клонов в $4$-значной логике}

\begin{document}
\maketitle
\tableofcontents
\newpage
\section{Перестановка переменных}
	Рассмотрим преобразование:
	\[f(x_{\psi(0)},x_{\psi(1)},\dots,x_{\psi(n)})\]
	Всегда можно так переставить переменные, чтобы полученная вектор-функция была лексико-графически минимальной.
	Также для любой лексико-графически минимальной вектор-функции существует ветка,
		в каждом узле которой вектор-функция являлась лексико-графически минимальной (при игнорировании нераскрытых значений).
	Это справедливо из-за свойства: если вектор-функция минимальна, то она была минимальной на всех узлах своейй ветви,
		ведь если её можно было уменьшить, то и итоговую функцию можно будет уменьшить.

	В таком случае можно отсекать те ветви, в которых вектор-функция не минимальна.

	Пусть мы добавляем значение в точке $\tilde a=(a_0,a_1,\dots a_n)$.
	Тогда (в худшем случае) нужно проверить все значения перестановок, в точках, лексикографически меньше, чем $\tilde a$.
	В некоторых случаях этого делать не нужно.
	Например, если при всех наборах веса $1$ разные значения, то уже не важно какие значения у остальных переменных.
\section{Перестановка значений}
	Рассмотрим преобразование сопряжения:
	\[f(\tilde x^n)\rightarrow\phi^{-1}f(\phi \tilde x^n)\]
	Это пребразование имеет смысл переименования элементов множества $E_k$.
	В таком смысле в таблице значений каждый элемент либо уже знаком, либо является новым.
	При рассмотрении функции с точностью до сопряжения, не важно какой это из новых элементов.
	Для определённости будем считать, что это минимальный неизвестный элемент.
\section{Общая фильтация}
	Рассмотрим частично заданнуую функцию $k$-значной логики $f(\tilde x)$.
	Для простоты пусть функция бинарная $f(\tilde x)=f(x,y)$.
	Тепперь рассмотрим $g(x,y) = f(x,f(x,y))$.
	Пусть она лексикографически меньше, чем $f(x,y)$.\\
	Если $f$ порождает минимальный клон,
		то $g$ тоже порождает минимальный клон.\\
	Действительно, если клон, порождённый $g$ не является минимальным,
		то клон, порождённый $f$ также не минимален.\\
	Раз функция $g$ появится в другой ветке (ведь $g<f$),
		то ветку с $f$ можно обрезать.

\section{Унарные функции}
	Для унарных функций $k$-значной логики, порождающих минимальный клон,
		теорема из статьи для $k=3$ легко обобщается на произвольные $k$.

	Нетривиальная унарная функция $f$ на конечном множестве порождает минимальный клон тогда и только тогда,
		когда $f$ является либо ретракцией $M$ (m.e. $f \circ f = f$), либо перестановкой простого порядка.
\section{Бинарные функции}
	Для бинарных функций $k$-значной логики, порождающих минимальный клон,
		теорема из статьи для $k=3$ также легко обобщается на произвольные $k$.

	Нетривиальная бинарная функция $f$ на конечном множестве является острой тогда и только тогда,
		когда она идемпотентна.
\end{document}
